\section*{Ensayo por Diego Alexander Cisneros Tenorio}
\addcontentsline{toc}{subsubsection}{Ensayo por Diego Alexander Cisneros Tenorio.}

Durante las últimas décadas, se han desarrollado tecnologías que han reestructurado por completo el estilo de vida de las personas; sin embargo, sus beneficios vienen con una serie de cuestiones éticas sobre la privacidad y la autonomía de los usuarios. Las tecnologías producen grandes cantidades de datos, un recurso intangible que permite obtener información sobre el comportamiento de los usuarios y su vida personal. Por esta razón, es importante tomar conciencia sobre la influencia y el poder que se les otorga a las corporaciones.

Así, surge una cuestión: ¿los usuarios deben ser los dueños de esta información? En lo personal, no creo que los usuarios tengamos la capacidad de realizar un uso responsable de todos nuestros datos personales; no obstante, la información producida por estas tecnologías es valiosa, por lo que, más allá de darles un control a los usuarios sobre cada parte de sus propios datos, creo que es mejor diseñar un estándar de normas éticas que regulen el uso y la información obtenida a partir de los datos recopilados.

Las redes sociales, el IOT y los servicios tecnológicos están tan presentes en la vida diaria que el no usarlos implica excluirse de la sociedad. De este modo, los usuarios permiten que las corporaciones recopilen y se beneficien con el valor de los datos generados por los usuarios. Por esta razón, el objetivo de la autora es convencer a los usuarios de tomar conciencia sobre el valor de sus datos y tener un control sobre ellos, de modo que las personas puedan deliberar y beneficiarse en caso de compartir sus datos con las grandes corporaciones, con el objetivo de limitar el poder de las corporaciones sobre los usuarios.

Personalmente, creo que la mayoría de las personas no están preparadas para tener y hacer uso de sus datos de forma responsable. Por ejemplo, en la mayoría de los casos, las personas suelen aceptar los términos y condiciones de los servicios digitales sin enterarse de los datos compartidos a las empresas. Por esta razón, considero que, antes de lograr que los usuarios se vuelvan conscientes del valor de sus datos, debemos prestar atención en regular la información que se les permite recopilar a las corporaciones.

Un obstáculo a la visión de la autora viene al considerar la enorme cantidad de datos que puede producir un usuario al usar un servicio. De acuerdo con el artículo, los usuarios de internet generan el 70 \% del universo digital y, en un escenario donde el tráfico de datos crece de forma exponencial cada año, implica que los usuarios deben hacerse cargo de más información cada vez. De este modo, resulta insostenible para los usuarios manejar y deliberar conscientemente qué parte de su información brindar a las corporaciones y qué parte mantenerla oculta. Además, según el artículo, los datos de forma individual no tienen un valor económico tan grande como los datos colectivos, por lo que resulta poco factible que los usuarios se beneficien económicamente al vender sus propios datos.

Por una parte, es cierto que una consecuencia de la recopilación de datos es la mejora en la experiencia de los usuarios. La recolección de datos sobre los usuarios proporciona una ventaja mutua entre el usuario y el proveedor. No obstante, sin medidas que regulen la recolección y el uso de datos sobre sus usuarios, las empresas pueden extraer información necesaria de los datos recopilados para crear perfiles particulares de sus usuarios y encontrar patrones en su comportamiento para atraerlos a sus productos, de forma que esta ventaja inicial puede transformarse en un medio de manipulación de masas.

Debemos considerar que la información obtenida por las corporaciones a partir de los servicios que brindan es una mina de información sobre comportamiento humano y es sumamente valiosa. Por esta razón, considero de suma importancia tomar medidas éticas que regulen los datos que pueden ser usados para mejorar la experiencia de los usuarios y qué datos son demasiado personales y no deberían de ser utilizados por las grandes corporaciones. Estoy de acuerdo en que no podemos esperar que las grandes corporaciones tomen medidas éticas sobre el uso de nuestros datos; por esta razón, considero que las normas éticas deben venir por parte de instituciones gubernamentales que regulen el uso de la información de los usuarios en beneficio de los usuarios. Promoviendo normas como las aplicadas por la GDPR de la Unión Europea, en las que se prohíbe recolectar datos sensibles y le otorgan al usuario la autoridad de exigir la eliminación de sus datos en cualquier momento.

Del mismo modo, en una sociedad sumergida en datos, es necesario mencionar el papel fundamental de los administradores de bases de datos, pues son los responsables de cumplir con las normas de seguridad al diseñar los sistemas de bases de datos. Por una parte, deben considerar aspectos de estructura, organización y de integridad en el sistema para garantizar que la información obtenida sea consistente y pueda ser usada en beneficio del cliente. Y, por otra parte, colocar especial atención en la seguridad, con el objetivo de proteger los datos almacenados y evitar que los datos de sus usuarios sean filtrados y sean consultados por entidades no autorizadas.

En conclusión, considero que es indispensable garantizar que los datos de los usuarios se almacenen de forma segura; sin embargo, otorgarles dicha responsabilidad a los usuarios se convierte más en una carga que en un beneficio a largo plazo. Reitero que establecer normas éticas sobre el uso de los datos y su correcto almacenamiento en una base de datos es la mejor forma de blindar nuestra información contra el uso abusivo que le puedan dar las corporaciones. Por otra parte, es notorio cómo las nuevas tecnologías se vuelven cada vez más cercanas a los individuos. Actualmente, con la llegada de la inteligencia artificial, las personas son capaces de soltar información más íntima, facilitándole el proceso a las corporaciones de recopilar datos y diseñar mejores perfiles sobre sus clientes, lo que supone un nuevo reto para establecer los límites éticos sobre el uso de la información.

En lo personal, este análisis me permite comprender la responsabilidad que recae en los diseñadores de bases de datos, reconociendo que su seguridad no solo es importante para la corporación, sino que, en un mundo de datos, también es fundamental para los usuarios.
